\exercisesection{}

\begin{enumerate}
\def\labelenumi{\arabic{enumi})}
\item
  Give an example of a group G, a field K, and a faithful representation of G such that the associated \(K[G]\)-module is not faithful.
\item
  Let \(\pi\) be a unitary representation of G on a Hilbert space E and \(x \in E\) of norm 1 such that
\end{enumerate}

\[
\sup _ {g \in \mathrm{G}} \| \pi (g) x - x \| \leqslant 1.
\]

There exists a nonzero G-invariant vector \(y \in E\).

\begin{enumerate}
\def\labelenumi{\arabic{enumi})}
\setcounter{enumi}{2}
\tightlist
\item
  Let E be a Hilbert space. We denote by \(\mathbf{U}(\mathbf{E})\) the group of unitary endomorphisms of E and \(\mathbf{GL}(\mathbf{E})\) the group of invertible endomorphisms of E.
\end{enumerate}

\begin{enumerate}
\def\labelenumi{\alph{enumi})}
\item
  The group \(\mathbf{U}(\mathrm{E})\) endowed with the topology induced by the topology of simple convergence on \(\mathcal{L}(\mathrm{E})\) is a topological group, denoted \(\mathbf{U}(\mathrm{E})_{s}\) ; if E is infinite-dimensional, the group \(\mathbf{GL}(\mathrm{E})\) is not a topological group with the topology induced by the topology of simple convergence.
\item
  The topology of \(\mathbf{U}(\mathrm{E})_{s}\) coincides with the topology induced by the topology of compact convergence.
\item
  If E is countably generated, the group \(\mathbf{U}(\mathrm{E})_{s}\) is metrizable and complete.
\item
  Let L be the space of \(u \in \mathcal{L}(\mathrm{E})\) such that \(u^{*} = -u\) ; it is the Lie algebra of the real Lie group G = U(E) endowed with the topology induced by the Banach space topology of \(\mathcal{L}(\mathrm{E})\) (LIE, III, p. 146, cor. 2). The exponential map from L to \(\mathbf{U}(\mathrm{E})_s\) is continuous.
\item
  If E is of infinite dimension, then \(\mathbf{U}(\mathrm{E})_{s}\) is not a real Lie group of type \(\{L\}\) . (Consider the subgroup K of \(\mathbf{U}(\mathrm{E})_{s}\) formed by the unitary endomorphisms diagonal in a fixed orthonormal basis of E, and note that K is compact.)
\item
  The exponential map is not a local homeomorphism from L to \(\mathbf{U}(\mathbf{E})_{s}\) . (Consider its restriction to the diagonal endomorphisms in a fixed orthonormal basis.)
\end{enumerate}

\begin{enumerate}
\def\labelenumi{\arabic{enumi})}
\setcounter{enumi}{3}
\tightlist
\item
  Let G be a topological group. We say that G is amenable if there exists a linear form \(\lambda\colon\mathcal{C}_{b}(\mathrm{G};\mathbf{R})\to\mathbf{R}\) such that
\end{enumerate}

\begin{enumerate}
\def\labelenumi{(\roman{enumi})}
\item
  If \(f \geqslant 0\) , then \(\lambda(f) \geqslant 0\) .
\item
  One has \(\lambda(1)=1\) .
\item
  For every \(g \in \mathbf{G}\) and every \(f \in \mathcal{C}_b(\mathbf{G};\mathbf{R})\) , we have \(\lambda(\boldsymbol{\gamma}(g)f) = \lambda(\boldsymbol{\delta}(g)f)\) , where
\end{enumerate}

\[
\boldsymbol {\gamma} (g) f (x) = f \left(g ^ {- 1} x\right), \quad \boldsymbol {\delta} (g) f (x) = f (x g).
\]

If G is a discrete group, this definition coincides with that of EVT, IV, p. 73, exerc. 4.

\begin{enumerate}
\def\labelenumi{\alph{enumi})}
\tightlist
\item
  The topological group G is amenable if and only if for every integer \(n \in N\) , every finite family \((f_i)_{1 \leqslant i \leqslant n}\) in \(\mathcal{C}_b(\mathbf{G}; \mathbf{R})\) , all families \((g_i)_{1 \leqslant i \leqslant n}\) and \((h_i)_{1 \leqslant i \leqslant n}\) in G, every real number a, the condition
\end{enumerate}

\[
\sum_ {i = 1} ^ {n} (f _ {i} (x) - f _ {i} (g _ {i} x h _ {i})) \geqslant a
\]

for all \(x \in G\) implies that \(a \leqslant 0\) . (To show that this condition is sufficient, consider the set of functions of the form

\[
x \mapsto \sum_ {i = 1} ^ {n} (f _ {i} (x) - f _ {i} (g _ {i} x h _ {i}))
\]

and apply the Hahn--Banach theorem.)

\begin{enumerate}
\def\labelenumi{\alph{enumi})}
\setcounter{enumi}{1}
\item
  Every compact topological group is amenable, and every abelian topological group is amenable.
\item
  Let G be a non-commutative free group. The discrete group G is not amenable.
\item
  For every Hilbert space E and every continuous bounded representation of G in E, there exists a Hilbert norm q on E such that the representation of G in the Hilbert space E endowed with q is unitary.
\end{enumerate}

\begin{enumerate}
\def\labelenumi{\arabic{enumi})}
\setcounter{enumi}{4}
\tightlist
\item
  Let G be a topological group and H an open subgroup of G. We denote by p the canonical projection \(G \rightarrow G/H\) .
\end{enumerate}

\begin{enumerate}
\def\labelenumi{\alph{enumi})}
\item
  Let \(\sigma\colon\mathrm{G/H}\to\mathrm{G}\) be a map such that \(p\circ\sigma=\mathrm{Id}_{\mathrm{G/H}}\) . We denote by \(\beta\colon\mathrm{G}\times\mathrm{G/H}\to\mathrm{H}\) the map defined by \(\beta(g,x)=\sigma(gx)^{-1}g\) . The map \(\beta\) is continuous.
\item
  Let \(\pi\) be a bounded representation of H in a Hilbert space E. We denote by F the Hilbert space \(\ell^2 (\mathrm{G / H};\mathrm{E})\) . The formula
\end{enumerate}

\[
\varrho (g) f (x) = \pi (\beta (g ^ {- 1}, x) ^ {- 1}) f (g ^ {- 1} x)
\]

defines a continuous bounded representation of G in F.

\begin{enumerate}
\def\labelenumi{\alph{enumi})}
\setcounter{enumi}{2}
\tightlist
\item
  Suppose there exists a Hilbert norm q on F such that the representation \(\varrho\) is a unitary representation of G in \((F,q)\) . Then there exists a Hilbert norm \(\widetilde{q}\) on E such that \(\pi\) is a unitary representation of H in \((E,\widetilde{q})\) .
\end{enumerate}

\begin{enumerate}
\def\labelenumi{\arabic{enumi})}
\setcounter{enumi}{5}
\tightlist
\item
  \begin{enumerate}
  \def\labelenumii{\alph{enumii})}
  \tightlist
  \item
    Let E be a complex Banach space and let \(u \in \mathcal{L}(\mathrm{E})\) . If there exists \(n \in N\) such that \(n \geqslant 1\) and \(u^{n} = 1_{E}\) then the spectrum of u is contained in the group of n-th roots of unity and consists of eigenvalues of u. b) Let G be a commutative topological group such that every element of G has finite order. Every continuous irreducible representation of G in a Banach space has dimension 1, and takes values in the group \(U \subset C^{*}\) of complex numbers of modulus 1.
  \end{enumerate}
\end{enumerate}

\begin{enumerate}
\def\labelenumi{\alph{enumi})}
\setcounter{enumi}{2}
\tightlist
\item
  Let G be a commutative topological group such that the set of elements of finite order of G is dense in G. Every continuous irreducible representation of G in a Banach space has dimension 1, and takes values in the group \(U \subset C^{*}\) of complex numbers of modulus 1.
\end{enumerate}

\begin{enumerate}
\def\labelenumi{\arabic{enumi})}
\setcounter{enumi}{6}
\tightlist
\item
  For every integer \(n \geqslant 1\) , we denote by \(s_{n}\) the function \(Z \rightarrow R\) defined by
\end{enumerate}

\[
s _ {n} (k) = \left\{ \begin{array}{l l} (1 - k) ^ {n (1 - k)} & \text{if } k \leqslant - 1 \\ 1 & \text{if } k = 0 \\ (k + 1) ^ {- (k + 1) / n} & \text{if } k \geqslant 1. \end{array} \right.
\]

We denote by S the set of functions \(s_n\) .

\begin{enumerate}
\def\labelenumi{\alph{enumi})}
\item
  For every integer \(n \geqslant 1\) and every \((k_{1}, k_{2}) \in \mathbf{Z}^{2}\) , we have \(s_{4n}(k_{1})s_{4n}(k_{2}) \geqslant s_{n}(k_{1} + k_{2})\) .
\item
  For \(f \in \mathbf{C}(\mathrm{T})\) , we denote by \((a_k(f))_{k \in \mathbf{Z}}\) the coefficients of the Laurent expansion of \(f\) at 0. For every \(s \in \mathrm{S}\) and every \(f \in \mathbf{C}(\mathrm{T})\) , the series
\end{enumerate}

\[
\sum_ {k \in \mathbf {Z}} | a _ {k} (f) | s (k)
\]

converges, and the sets

\[
\mathcal {V} (s, \varepsilon) = \{f \in \mathbf {C} (\mathrm{T}) | \sum_ {k \in \mathbf {Z}} | a _ {k} (f) | s (k) <   \varepsilon \}
\]

defined for \(s \in S\) and \(\varepsilon > 0\) , form a fundamental system of open neighborhoods of 0 for a locally convex topology on \(\mathbf{C}(T)\) . The multiplication \(\mathbf{C}(T) \times \mathbf{C}(T) \to \mathbf{C}(T)\) is continuous for this topology.

\begin{enumerate}
\def\labelenumi{\alph{enumi})}
\setcounter{enumi}{2}
\tightlist
\item
  The topology defined in the preceding question is metrizable but not complete.
\end{enumerate}

\(d)\) For \(f \in \mathbf{C}(\mathrm{T})^*\) , we denote by \(\varrho(f)\) the map \(g \mapsto fg\) from \(\mathbf{C}(\mathrm{T})\) into \(\mathbf{C}(\mathrm{T})\) . The map \(\varrho\) is a continuous representation of \(\mathbf{C}(\mathrm{T})^*\) in the locally convex space \(\mathbf{C}(\mathrm{T})\) .

\begin{enumerate}
\def\labelenumi{\alph{enumi})}
\setcounter{enumi}{4}
\tightlist
\item
  The representation \(\varrho\) is irreducible.
\end{enumerate}

